% PRINCIPIA PHYSICA -- pilot topic 2
% Hamiltonian reconstruction as a physical theory object
% arXiv-style preprint, self-contained.
\documentclass[10pt,a4paper]{article}

\usepackage[margin=0.9in]{geometry}
\usepackage[T1]{fontenc}
\usepackage{amsmath,amssymb,amsthm,mathtools}
\usepackage{microtype}
\usepackage{enumitem}
\usepackage{listings}
\usepackage{multicol}
\usepackage[compact]{titlesec}
\titlespacing*{\section}{0pt}{6pt}{3pt}
\titlespacing*{\subsection}{0pt}{4pt}{2pt}
\usepackage[colorlinks=true,linkcolor=blue,citecolor=blue,urlcolor=blue]{hyperref}

\makeatletter
\def\thm@space@setup{\thm@preskip=3pt \thm@postskip=3pt}
\makeatother

\theoremstyle{plain}
\newtheorem{theorem}{Theorem}
\newtheorem{proposition}[theorem]{Proposition}
\newtheorem{lemma}[theorem]{Lemma}
\newtheorem{corollary}[theorem]{Corollary}

\theoremstyle{definition}
\newtheorem{definition}[theorem]{Definition}
\newtheorem{axiom}[theorem]{Axiom}
\newtheorem{assumption}[theorem]{Assumption}
\newtheorem{conjecture}[theorem]{Conjecture}
\newtheorem{interpretation}[theorem]{Interpretation}
\newtheorem{empirical}[theorem]{Empirical statement}
\newtheorem{established}[theorem]{Established physical result}
\newtheorem{newresult}[theorem]{New result claimed by this work}
\newtheorem{example}[theorem]{Example}
\newtheorem{nonexample}[theorem]{Non-example}
\newtheorem{remark}[theorem]{Remark}

\newcommand{\Obs}{\mathrm{Obs}}
\newcommand{\Ham}{\mathbf{Ham}}
\newcommand{\HamP}{\mathbf{Ham}_{\mathrm{port}}}
\newcommand{\Cinf}{C^{\infty}}
\newcommand{\dR}{\mathrm{dR}}
\newcommand{\Lie}{\mathcal{L}}
\newcommand{\ins}{\iota}
\newcommand{\R}{\mathbb{R}}
\newcommand{\Zed}{\mathbb{Z}}
\newcommand{\PB}[2]{\{#1,#2\}}
\DeclareMathOperator{\Jac}{Jac}
\DeclareMathOperator{\tr}{tr}
\DeclareMathOperator{\Sym}{Sym}

\lstset{
  basicstyle=\ttfamily\scriptsize,
  aboveskip=4pt, belowskip=4pt,
  columns=fullflexible,
  keepspaces=true,
  frame=single,
  breaklines=true,
  showstringspaces=false
}

\title{\bf Hamiltonian Reconstruction as a Physical Theory Object\\
\large Finite-dimensional classical mechanics as a typed, composable, empirically interpreted structure}

\author{PRINCIPIA PHYSICA Pilot, Topic 2\\
\small Autonomous research worker, AgentHero research market}

\date{15 August 2026}

\begin{document}
\maketitle

\begin{abstract}
We reconstruct finite-dimensional Hamiltonian mechanics as a \emph{typed theory object}: a four-layer structure consisting of a kinematic layer (a symplectic manifold), a dynamical layer (a distinguished generator), an interface layer (designated port observables), and an empirical interpretation layer (dimensioned measurement procedures with stated uncertainties), arranged so that the mathematical layers are derived from explicit axioms while the empirical layer is never derived from them. We prove: (i) a reconstruction theorem, that the Hamiltonian form of the generator is forced by structure preservation plus vanishing first de~Rham cohomology, with an explicit torus counterexample when that cohomology is nonzero; (ii) a compositional theorem, that a linearly coupled pair of Hamiltonian objects preserves the \emph{product} symplectic form iff the coupling is reciprocal, and, for one degree of freedom per factor, preserves \emph{some} split-respecting symplectic form iff the two coupling constants vanish together, making unidirectional coupling a sharp obstruction to composition; (iii) a non-conflation theorem, that isomorphism of the closed-system layers neither implies nor is implied by empirical equivalence, witnessed by an isomorphic-yet-separable oscillator pair and a non-isomorphic-yet-inseparable pendulum pair. We give non-examples (damped oscillator, non-Jacobi bracket, $S^6$, unidirectional coupling), scope limits, falsification conditions, and a machine-executed Haskell appendix checking the composition theorems. Every claim carries an explicit epistemic label. We do not claim that nature is Hamiltonian, only that a precisely stated set of assumptions forces the Hamiltonian form, and we exhibit what would falsify each assumption.
\end{abstract}

\section{Introduction}
\label{sec:intro}

\subsection{The thesis under test}

The PRINCIPIA PHYSICA programme investigates the thesis that \emph{physics is the study of empirically realized compositional mathematical structures}. This paper does not assume that thesis. It performs one bounded test: it takes the best-developed piece of mathematical physics in the finite-dimensional regime---Hamiltonian mechanics---and asks whether it can be presented as a \emph{single typed object} whose mathematical content is generated by stated axioms, whose empirical content is carried by a separate and irreducible layer, and whose composition operation has an explicitly characterised domain of validity. A clean success would be weak confirmation that the thesis is well-posed here; success purchased by smuggling empirical content into the mathematical layer would be evidence against it; failure at composition would show that ``compositional'' does less work than it appears. We find an intermediate case: composition succeeds, but only on a proper subclass of couplings that we characterise exactly (Section~\ref{sec:composition}), and the obstruction is physically realized, not a technicality.

\begin{assumption}[Scope]
\label{ass:scope}
Throughout, all manifolds are smooth ($\Cinf$), second countable, Hausdorff, and finite dimensional; all maps are smooth; all vector fields are complete unless stated otherwise. Field theory, infinite-dimensional analysis, general relativity, quantum mechanics, quantum field theory, and the geometry of theory space are outside scope and appear only as boundary tests in Section~\ref{sec:limits}. No new experimental claim is made anywhere in this paper.
\end{assumption}

\subsection{Contributions, epistemic status, and conventions}

Definition~\ref{def:object} (typed signature) and Axiom~\ref{ax:sep} (layer separation) are definitions, not discoveries; their value is that they make the later negative results statable. Theorem~\ref{thm:reconstruction} assembles classical facts \cite{AbrahamMarsden1978,Arnold1989,MarsdenRatiu1999}, stated in typed form with proof in full because the exact hypotheses matter for the falsification analysis of Section~\ref{sec:falsification}. Theorems~\ref{thm:reciprocity} and \ref{thm:splitobstruction} are claimed new in this precise form, though the first is a coordinate restatement of standard reciprocity and the second is elementary linear algebra; no working physicist would be surprised by either. Theorem~\ref{thm:noniso} and Proposition~\ref{prop:nonnecessary} formalise and prove, in a specific category, the rule that mathematical equivalence is not physical equivalence; the physics inside those proofs (mass is measurable; a pendulum is nearly harmonic) is elementary and long known, and only the framework-level statement is claimed.

Every numbered environment carries one of the labels: Definition, Axiom, Assumption, Conjecture, Lemma, Proposition, Theorem, Corollary, Interpretation, Empirical statement, Established physical result, New result claimed by this work. Assumptions are never upgraded to theorems by later use; statements labelled Established physical result are attributed and not re-derived; statements labelled Empirical statement report measurement content and are not proved. We avoid ``clearly'', ``obviously'', ``it follows'', and ``one can show'' as substitutes for argument.

\paragraph{Sign conventions.} We fix, once, the convention
\begin{equation}
\label{eq:conventions}
\ins_{X_f}\omega = df, \qquad \PB{f}{g} := \omega(X_f,X_g),
\end{equation}
so that in canonical coordinates with $\omega = \sum_i dq^i\wedge dp_i$ we obtain $X_H = \sum_i\big(\partial_{p_i}H\,\partial_{q^i} - \partial_{q^i}H\,\partial_{p_i}\big)$, Hamilton's equations $\dot q^i = \partial H/\partial p_i$, $\dot p_i = -\partial H/\partial q^i$, the classical bracket $\PB{f}{g} = \sum_i(\partial_{q^i}f\,\partial_{p_i}g - \partial_{p_i}f\,\partial_{q^i}g)$, and the identity $X_g f = \PB{f}{g}$. With this convention $X_{\PB{f}{g}} = -[X_f,X_g]$; the opposite sign appears in texts using $\ins_{X_f}\omega = -df$. Convention dependence of this sign is discussed in \cite[Ch.~3]{AbrahamMarsden1978}.

\section{Typed theory objects}
\label{sec:typed}

\begin{definition}[Type discipline]
\label{def:types}
A \emph{typed presentation} of a physical theory is a presentation in which every symbol carries a declared type, every judgement carries one of the epistemic labels listed above, and no derivation rule permits a judgement of label $\ell$ to be inferred from premises unless the inference is licensed for $\ell$. In particular: an Assumption may be used as a premise but never appears as a conclusion of a proof; an Empirical statement may not be the conclusion of a purely mathematical derivation.
\end{definition}

\begin{definition}[Physical theory object]
\label{def:object}
A \emph{physical theory object} is a tuple
\[
\mathfrak{T} \;=\; \big(\, \mathcal{K},\ \mathcal{D},\ \mathcal{I},\ \mathcal{E} \,\big)
\]
of four layers:
\begin{description}[leftmargin=1.4em,itemsep=1pt]
\item[$\mathcal{K}$ (kinematic layer)] a set of states with distinguished structure, and a commutative $\R$-algebra $\Obs$ of observables acting on it;
\item[$\mathcal{D}$ (dynamical layer)] a distinguished one-parameter family of endomorphisms of $\mathcal{K}$ preserving the distinguished structure, together with its infinitesimal generator;
\item[$\mathcal{I}$ (interface layer)] a finite tuple $(A_1,\dots,A_k)$ of elements of $\Obs$, the \emph{ports}, through which the object may be composed with other objects;
\item[$\mathcal{E}$ (empirical interpretation layer)] a triple $(P,\rho,\varepsilon)$ where $P$ is a finite set of laboratory procedures, $\rho: P \to \Obs \times D$ assigns to each procedure an observable together with a dimension in a finitely generated free abelian dimension group $D$, and $\varepsilon: P \to \R_{>0}$ assigns a stated uncertainty.
\end{description}
A \emph{morphism} of theory objects preserves all four layers; forgetting layers gives forgetful functors between the resulting categories, in the sense of \cite[Ch.~IV]{MacLane1998}.
\end{definition}

\begin{axiom}[Layer separation]
\label{ax:sep}
$\mathcal{E}$ is not a function of $(\mathcal{K},\mathcal{D},\mathcal{I})$. No theorem of this paper derives any component of $\mathcal{E}$ from the other three layers, and no assignment of $\mathcal{E}$ is claimed to be canonical.
\end{axiom}

Axiom~\ref{ax:sep} is methodological, adopted so that the rule ``mathematical equivalence must not be identified with physical equivalence without an empirical argument'' is enforceable rather than merely asserted. Without it, Theorem~\ref{thm:noniso} could not be stated.

\begin{interpretation}
\label{int:whyfour}
The layers answer four different questions: \emph{what can be the case} ($\mathcal{K}$), \emph{what happens} ($\mathcal{D}$), \emph{what can be attached} ($\mathcal{I}$), \emph{what is read off an instrument} ($\mathcal{E}$). A structure lacking $\mathcal{E}$ is mathematics; one lacking $\mathcal{I}$ is a closed universe and cannot be experimented on; one lacking $\mathcal{D}$ is kinematics. That all four are needed is an interpretive claim of this paper, not a theorem.
\end{interpretation}

\begin{example}[A fully specified object]
\label{ex:osc}
Let $\mathfrak{h}_{\mathrm{osc}}$ have $\mathcal{K} = \big(\R^2,\ dq\wedge dp,\ \Obs = \Cinf(\R^2)\big)$; $\mathcal{D}$ generated by $H = \frac{p^2}{2m} + \frac12 m\omega_0^2q^2$; $\mathcal{I} = (q)$; and $\mathcal{E} = (P,\rho,\varepsilon)$ with $P = \{\mathsf{disp},\mathsf{per}\}$, $\rho(\mathsf{disp}) = (q,\ \mathrm{L})$ read by a calibrated interferometer, $\rho(\mathsf{per}) = (2\pi/\omega_0,\ \mathrm{T})$ read by a clock, dimensions valued in the free abelian group on $\{\mathrm{L},\mathrm{M},\mathrm{T}\}$, and $\varepsilon$ the instruments' stated uncertainties. All four layers are specified independently; recalibrating the interferometer changes $\mathcal{E}$ alone and yields a different object with identical $(\mathcal{K},\mathcal{D},\mathcal{I})$. This object recurs in Theorem~\ref{thm:noniso}.
\end{example}

\section{The kinematic layer}
\label{sec:kinematic}

\begin{definition}[Symplectic manifold]
\label{def:symp}
A \emph{symplectic manifold} is a pair $(M,\omega)$ with $M$ a smooth manifold and $\omega \in \Omega^2(M)$ satisfying (a) \emph{nondegeneracy}: for each $x \in M$ the map $\omega^\flat_x: T_xM \to T^*_xM$, $v \mapsto \omega_x(v,\cdot)$, is a linear isomorphism; and (b) \emph{closedness}: $d\omega = 0$. Nondegeneracy forces $\dim M = 2n$ even. We take $\Obs := \Cinf(M,\R)$. Standard references for this layer are \cite{AbrahamMarsden1978,Arnold1989,LibermannMarle1987,McDuffSalamon2017}.
\end{definition}

\begin{lemma}[Existence and uniqueness of the Hamiltonian vector field]
\label{lem:XH}
For each $H \in \Cinf(M)$ there is exactly one vector field $X_H$ with $\ins_{X_H}\omega = dH$.
\end{lemma}
\begin{proof}
By nondegeneracy $\omega^\flat_x$ is an isomorphism at each $x$, so $X_H(x) := (\omega^\flat_x)^{-1}(dH_x)$ is the unique pointwise solution. Since $\omega$ and $dH$ are smooth and $\omega^\flat$ is a smooth bundle isomorphism, its inverse is smooth, hence $X_H$ is smooth. If $\ins_{X}\omega = \ins_{X'}\omega$ then $\omega^\flat(X - X') = 0$, so $X = X'$.
\end{proof}

\begin{lemma}[Pointwise surjectivity]
\label{lem:surj}
For every $x \in M$ and every $v \in T_xM$ there is $f \in \Cinf(M)$ with $X_f(x) = v$.
\end{lemma}
\begin{proof}
Let $\alpha := \omega^\flat_x(v) \in T^*_xM$. In a chart around $x$ with coordinates $(u^1,\dots,u^{2n})$ centred at $x$, put $\tilde f := \sum_i \alpha_i u^i$ and let $f := \chi \tilde f$ with $\chi$ a bump function equal to $1$ near $x$ and compactly supported in the chart; such $\chi$ exists by \cite[Ch.~2]{Lee2013}. Then $df_x = \alpha$, hence $X_f(x) = (\omega^\flat_x)^{-1}\alpha = v$.
\end{proof}

\begin{definition}[Bracket]
\label{def:bracket}
For $f,g \in \Cinf(M)$ set $\PB{f}{g} := \omega(X_f,X_g)$. It is $\R$-bilinear, antisymmetric, and a derivation in each argument, the last because $X_{fg} = f X_g + g X_f$ follows from $d(fg) = f\,dg + g\,df$ and Lemma~\ref{lem:XH}.
\end{definition}

The next proposition is the precise sense in which the closedness axiom of Definition~\ref{def:symp} is \emph{not} a technical convenience: it is exactly the Jacobi identity, hence exactly Poisson's theorem on constants of motion.

\begin{proposition}[Closedness equals Jacobi]
\label{prop:jacobi}
Let $\omega$ be a nondegenerate $2$-form on $M$ and define $X_f$ and $\PB{\cdot}{\cdot}$ by \eqref{eq:conventions}. Write $\Jac(f,g,h) := \PB{\PB{f}{g}}{h} + \PB{\PB{g}{h}}{f} + \PB{\PB{h}{f}}{g}$. Then for all $f,g,h \in \Cinf(M)$,
\begin{equation}
\label{eq:jacobi-identity}
d\omega(X_f,X_g,X_h) \;=\; 2\,\Jac(f,g,h).
\end{equation}
Consequently $d\omega = 0$ if and only if $\PB{\cdot}{\cdot}$ satisfies the Jacobi identity.
\end{proposition}
\begin{proof}
Use the intrinsic formula for the exterior derivative of a $2$-form,
\begin{multline*}
d\omega(X,Y,Z) = X\,\omega(Y,Z) - Y\,\omega(X,Z) + Z\,\omega(X,Y)\\
{} - \omega([X,Y],Z) + \omega([X,Z],Y) - \omega([Y,Z],X),
\end{multline*}
valid for any vector fields \cite[Ch.~14]{Lee2013}, with $X = X_f$, $Y = X_g$, $Z = X_h$. Using $\omega(X_a,X_b) = \PB{a}{b}$ and $X_a b = \PB{b}{a}$, the first three terms are
\[
\PB{\PB{g}{h}}{f} - \PB{\PB{f}{h}}{g} + \PB{\PB{f}{g}}{h}.
\]
For the last three, note that $X_{\PB{a}{b}} = -[X_a,X_b]$; this is the standard identity for the convention \eqref{eq:conventions} \cite[Prop.~3.3.17]{AbrahamMarsden1978}, and it is verified directly on $f = q^2/2$, $g = p^2/2$ in one degree of freedom, where $X_f = -q\,\partial_p$, $X_g = p\,\partial_q$, $[X_f,X_g] = -q\,\partial_q + p\,\partial_p$, $\PB{f}{g} = qp$ and $X_{qp} = q\,\partial_q - p\,\partial_p$. Hence
\[
-\omega([X_f,X_g],X_h) = \PB{\PB{f}{g}}{h}, \quad
+\omega([X_f,X_h],X_g) = -\PB{\PB{f}{h}}{g}, \quad
-\omega([X_g,X_h],X_f) = \PB{\PB{g}{h}}{f}.
\]
Summing the six terms and using $-\PB{\PB{f}{h}}{g} = \PB{\PB{h}{f}}{g}$ gives \eqref{eq:jacobi-identity}.

For the consequence: if $d\omega = 0$ then $\Jac \equiv 0$ by \eqref{eq:jacobi-identity}. Conversely if $\Jac \equiv 0$, fix $x$ and $u,v,w \in T_xM$; by Lemma~\ref{lem:surj} choose $f,g,h$ with $X_f(x) = u$, $X_g(x) = v$, $X_h(x) = w$; then $d\omega_x(u,v,w) = 2\Jac(f,g,h)(x) = 0$. As $x,u,v,w$ were arbitrary, $d\omega = 0$.
\end{proof}

\begin{nonexample}[A bracket without Jacobi]
\label{nonex:jacobi}
On $\R^3$ with coordinates $(x_1,x_2,x_3)$ let $\Pi := x_1\,\partial_1\wedge\partial_2 + x_3\,\partial_2\wedge\partial_3 + x_1\,\partial_3\wedge\partial_1$ and $\PB{f}{g} := \Pi(df,dg)$. This bracket is bilinear, antisymmetric and a derivation in each slot, with $\PB{x_1}{x_2} = x_1$, $\PB{x_2}{x_3} = x_3$, $\PB{x_3}{x_1} = x_1$. Then
\begin{align*}
\Jac(x_1,x_2,x_3) &= \PB{\PB{x_1}{x_2}}{x_3} + \PB{\PB{x_2}{x_3}}{x_1} + \PB{\PB{x_3}{x_1}}{x_2}\\
&= \PB{x_1}{x_3} + \PB{x_3}{x_1} + \PB{x_1}{x_2} = -x_1 + x_1 + x_1 = x_1 \not\equiv 0 .
\end{align*}
\end{nonexample}

\begin{proposition}[What Jacobi buys]
\label{prop:poissonthm}
Let $\PB{\cdot}{\cdot}$ be bilinear, antisymmetric, and a derivation in each argument on $\Cinf(M)$, and for $H \in \Cinf(M)$ let $X_H := \PB{\cdot}{H}$ act on observables. Then $X_H$ is a derivation of the bracket, i.e.\ $X_H\PB{f}{g} = \PB{X_H f}{g} + \PB{f}{X_H g}$ for all $f,g$, if and only if $\Jac(f,g,H) = 0$ for all $f,g$. In particular, if Jacobi holds and $f,g$ are constants of motion for $H$, so is $\PB{f}{g}$; and in Non-example~\ref{nonex:jacobi} this conclusion is unavailable.
\end{proposition}
\begin{proof}
Expand: $X_H\PB{f}{g} - \PB{X_Hf}{g} - \PB{f}{X_Hg} = \PB{\PB{f}{g}}{H} - \PB{\PB{f}{H}}{g} - \PB{f}{\PB{g}{H}}$. Antisymmetry converts the last two terms into $+\PB{\PB{h}{f}}{g}$-type terms with $h = H$: indeed $-\PB{\PB{f}{H}}{g} = \PB{\PB{H}{f}}{g}$ and $-\PB{f}{\PB{g}{H}} = \PB{\PB{g}{H}}{f}$. The sum is exactly $\Jac(f,g,H)$. The final assertion: if $\PB{f}{H} = \PB{g}{H} = 0$ then $\PB{\PB{f}{g}}{H} = X_H\PB{f}{g} = \PB{X_Hf}{g} + \PB{f}{X_Hg} = 0$.
\end{proof}

\begin{established}[Darboux]
\label{est:darboux}
Every symplectic manifold $(M^{2n},\omega)$ admits, around each point, coordinates $(q^1,\dots,q^n,p_1,\dots,p_n)$ in which $\omega = \sum_i dq^i\wedge dp_i$ \cite[Thm.~3.2.2]{AbrahamMarsden1978}, \cite[Ch.~8]{Arnold1989}, \cite[Thm.~3.2.2]{McDuffSalamon2017}. Hence the kinematic layer has no local invariants beyond dimension; all local structure resides in $\mathcal{D}$, $\mathcal{I}$, $\mathcal{E}$.
\end{established}

\begin{interpretation}
\label{int:darboux}
Established result~\ref{est:darboux} is why the interface and empirical layers cannot be dispensed with: if the kinematic layer alone fixed physical content, all $2n$-dimensional classical systems would be locally physically identical, which no laboratory reports. Darboux's theorem is thus indirect internal evidence for Axiom~\ref{ax:sep}, not a proof of it.
\end{interpretation}

\section{Reconstruction of the dynamical layer}
\label{sec:reconstruction}

We now state the assumptions from which the Hamiltonian form of the dynamics is derived, and derive it. The point of stating them separately is that each is independently falsifiable; Section~\ref{sec:falsification} says how.

\begin{assumption}[Dynamical axioms]
\label{ass:dyn}
Let $(M,\omega)$ be symplectic. Assume:
\begin{description}[leftmargin=2.2em,itemsep=0pt]
\item[(D1)] \emph{Determinism and reversibility}: the dynamics is a smooth action $\varphi: \R\times M \to M$ of the additive group $\R$, i.e.\ $\varphi_0 = \mathrm{id}$ and $\varphi_{s+t} = \varphi_s\circ\varphi_t$.
\item[(D2)] \emph{Structure preservation}: $\varphi_t^*\omega = \omega$ for all $t$.
\item[(D3)] \emph{Autonomy}: $\varphi$ does not depend on an external time label beyond (D1).
\end{description}
\end{assumption}

\begin{theorem}[Reconstruction]
\label{thm:reconstruction}
Under Assumption~\ref{ass:dyn}, let $X$ be the infinitesimal generator of $\varphi$. Then:
\begin{enumerate}[label=(\roman*),leftmargin=2em,itemsep=0pt]
\item $\ins_X\omega$ is a closed $1$-form;
\item every point of $M$ has a neighbourhood $U$ and $H \in \Cinf(U)$ with $X|_U = X_H$;
\item if $H^1_{\dR}(M) = 0$, there is a global $H \in \Cinf(M)$ with $X = X_H$, unique up to an additive constant on each connected component.
\end{enumerate}
\end{theorem}
\begin{proof}
(i) Differentiating (D2) at $t$ and using the definition of the Lie derivative, $0 = \frac{d}{dt}\varphi_t^*\omega = \varphi_t^*(\Lie_X\omega)$; since $\varphi_t$ is a diffeomorphism, $\Lie_X\omega = 0$. By Cartan's magic formula, $\Lie_X\omega = d\ins_X\omega + \ins_X d\omega$, and $d\omega = 0$ by Definition~\ref{def:symp}; hence $d(\ins_X\omega) = 0$.

(ii) By the Poincar\'e lemma \cite[Ch.~17]{Lee2013}, on any contractible open $U$ the closed $1$-form $\ins_X\omega$ is exact: $\ins_X\omega|_U = dH$ for some $H \in \Cinf(U)$. By Lemma~\ref{lem:XH}, $X|_U = X_H$.

(iii) If $H^1_{\dR}(M) = 0$ then the closed form $\ins_X\omega$ is globally exact, $\ins_X\omega = dH$, and $X = X_H$ by Lemma~\ref{lem:XH}. If also $X = X_{H'}$ then $d(H - H') = 0$, so $H - H'$ is locally constant, hence constant on connected components.
\end{proof}

\begin{remark}
\label{rem:nottheorem}
Theorem~\ref{thm:reconstruction} does not state that physical dynamics is Hamiltonian. It states that (D1)--(D3) plus symplecticity force the Hamiltonian form. Assumption~\ref{ass:dyn} remains an assumption wherever it is used below; it is not converted into a theorem by Theorem~\ref{thm:reconstruction}.
\end{remark}

\begin{nonexample}[Locally but not globally Hamiltonian]
\label{nonex:torus}
Let $M = \mathbb{T}^2 = \R^2/\Zed^2$ with angle coordinates $(\theta_1,\theta_2)$ and $\omega = d\theta_1\wedge d\theta_2$, which is nondegenerate and closed (it is a top form), so $(M,\omega)$ is symplectic and compact. Let $X = \partial_{\theta_1}$, whose flow $\varphi_t(\theta_1,\theta_2) = (\theta_1+t,\theta_2)$ satisfies (D1)--(D3). Then $\ins_X\omega = d\theta_2$, which is closed but not exact, since $\int_{\gamma} d\theta_2 = 1 \neq 0$ for the loop $\gamma(s) = (0,s)$. Hence $X$ is locally Hamiltonian but admits no global $H$ with $X = X_H$. Here $H^1_{\dR}(\mathbb{T}^2) \cong \R^2 \neq 0$, so hypothesis (iii) of Theorem~\ref{thm:reconstruction} fails, and the failure is exactly detected by the topology.
\end{nonexample}

\begin{corollary}[Liouville]
\label{cor:liouville}
Under Assumption~\ref{ass:dyn}, $\varphi_t$ preserves the volume form $\omega^n/n!$.
\end{corollary}
\begin{proof}
$\varphi_t^*(\omega^n) = (\varphi_t^*\omega)^n = \omega^n$ because pullback is an algebra homomorphism for the wedge product.
\end{proof}

The converse of Corollary~\ref{cor:liouville} fails for $2n \geq 4$; the witness is constructed in Section~\ref{sec:composition} (Remark~\ref{rem:volumeweak}). This matters for falsification: measuring volume preservation is \emph{not} the same as measuring Hamiltonicity.

\section{Composition, and the sharp obstruction to it}
\label{sec:composition}

\begin{definition}[Free composition]
\label{def:tensor}
Given theory objects $\mathfrak{h}_i = ((M_i,\omega_i,\Obs_i),\,H_i,\,(A^{(i)}_1,\dots),\,\mathcal{E}_i)$ for $i=1,2$, their \emph{free composite} is
\[
\mathfrak{h}_1\otimes\mathfrak{h}_2 := \big( (M_1\times M_2,\ \pi_1^*\omega_1 + \pi_2^*\omega_2),\ \ H_1\circ\pi_1 + H_2\circ\pi_2,\ \ \text{ports of both},\ \ \mathcal{E}_1 \sqcup \mathcal{E}_2 \big).
\]
Write $\omega_\oplus := \pi_1^*\omega_1 + \pi_2^*\omega_2$.
\end{definition}

\begin{proposition}[The free composite is a theory object and the parts remain identifiable]
\label{prop:tensorok}
$\omega_\oplus$ is symplectic on $M_1\times M_2$; each $\pi_i$ is a Poisson map; and $\PB{f\circ\pi_1}{g\circ\pi_2}_{\omega_\oplus} = 0$ for all $f \in \Cinf(M_1)$, $g \in \Cinf(M_2)$.
\end{proposition}
\begin{proof}
$d\omega_\oplus = \pi_1^*d\omega_1 + \pi_2^*d\omega_2 = 0$. In the splitting $T_{(x_1,x_2)}(M_1\times M_2) = T_{x_1}M_1\oplus T_{x_2}M_2$ we have $\omega_\oplus^\flat = \omega_1^\flat\oplus\omega_2^\flat$, a direct sum of isomorphisms, so $\omega_\oplus$ is nondegenerate. Since $d(f\circ\pi_1) = \pi_1^*df$ annihilates $TM_2$, the field $X_{f\circ\pi_1}$ is tangent to the $M_1$ factor, and likewise $X_{g\circ\pi_2}$ to $M_2$; as $\omega_\oplus$ pairs the factors trivially, $\PB{f\circ\pi_1}{g\circ\pi_2} = 0$. The same tangency computation shows each $\pi_i$ is Poisson.
\end{proof}

Free composition is not interaction. Interaction is where the compositional claim in the founding thesis is actually tested.

\begin{definition}[Linear position--position interconnection]
\label{def:coupling}
Let $M_i = T^*\R^{n_i} = \R^{2n_i}$ with $\omega_i = \sum_a dq_a^{(i)}\wedge dp^{(i)}_a$ and quadratic Hamiltonians giving the free equations $\dot q^{(i)} = M_i^{-1}p^{(i)}$, $\dot p^{(i)} = -K_i q^{(i)}$. An \emph{interconnection with coupling matrices} $C_{12} \in \R^{n_1\times n_2}$, $C_{21} \in \R^{n_2\times n_1}$ is the vector field $X$ on $M_1\times M_2$ given by
\begin{equation}
\label{eq:coupled}
\dot q^{(1)} = M_1^{-1}p^{(1)},\quad
\dot p^{(1)} = -K_1q^{(1)} - C_{12}q^{(2)},\quad
\dot q^{(2)} = M_2^{-1}p^{(2)},\quad
\dot p^{(2)} = -K_2q^{(2)} - C_{21}q^{(1)}.
\end{equation}
The ports are the position observables $q^{(1)}, q^{(2)}$.
\end{definition}

\begin{theorem}[Reciprocity is exactly strict compositionality]
\label{thm:reciprocity}
The flow of \eqref{eq:coupled} preserves $\omega_\oplus$ if and only if $C_{21} = C_{12}^{\!\top}$. In that case the composite is the Hamiltonian object with $H = H_1 + H_2 + H_{\mathrm{int}}$ and $H_{\mathrm{int}} = (q^{(1)})^{\!\top}C_{12}\,q^{(2)}$.
\end{theorem}
\begin{proof}
Write $X = X_{H_1+H_2} + Y$ with $Y = -\big(C_{12}q^{(2)}\big)\cdot\partial_{p^{(1)}} - \big(C_{21}q^{(1)}\big)\cdot\partial_{p^{(2)}}$. Since $\Lie_{X_{H_1+H_2}}\omega_\oplus = 0$ and $\Lie$ is linear in the field, $\Lie_X\omega_\oplus = \Lie_Y\omega_\oplus = d\,\ins_Y\omega_\oplus$ by Cartan and $d\omega_\oplus = 0$. For a field with $\partial_{p}$-components $b$ and no $\partial_q$-component, $\ins_Y(\sum dq\wedge dp) = -\,b\cdot dq$, so
\[
\ins_Y\omega_\oplus = \big(C_{12}q^{(2)}\big)\cdot dq^{(1)} + \big(C_{21}q^{(1)}\big)\cdot dq^{(2)}
= \sum_{a,b}(C_{12})_{ab}\,q^{(2)}_b\,dq^{(1)}_a + \sum_{a,b}(C_{21})_{ba}\,q^{(1)}_a\,dq^{(2)}_b .
\]
Hence $d\,\ins_Y\omega_\oplus = \sum_{a,b}\big((C_{21})_{ba} - (C_{12})_{ab}\big)\,dq^{(1)}_a\wedge dq^{(2)}_b$, which vanishes if and only if $(C_{21})_{ba} = (C_{12})_{ab}$ for all $a,b$, i.e.\ $C_{21} = C_{12}^{\top}$. When it vanishes, $\ins_Y\omega_\oplus = dH_{\mathrm{int}}$ with $H_{\mathrm{int}} = (q^{(1)})^{\top}C_{12}q^{(2)}$, so $Y = X_{H_{\mathrm{int}}}$ by Lemma~\ref{lem:XH}, and $X = X_{H_1+H_2+H_{\mathrm{int}}}$.
\end{proof}

\begin{remark}[Volume preservation is strictly weaker]
\label{rem:volumeweak}
The field \eqref{eq:coupled} is linear with $\tr A = 0$ for \emph{all} $C_{12},C_{21}$, since each $\dot q$ block is independent of $q$ and each $\dot p$ block of $p$. Its flow therefore preserves Lebesgue volume even when $C_{12}\neq C_{21}^{\top}$, where by Theorem~\ref{thm:reciprocity} it does not preserve $\omega_\oplus$: the promised witness that the converse of Corollary~\ref{cor:liouville} fails for $2n\geq 4$. The distinction is classical \cite[Ch.~8]{Arnold1989}.
\end{remark}

The natural objection to Theorem~\ref{thm:reciprocity} is that a non-reciprocal composite might still be Hamiltonian for \emph{some other} symplectic form. It might. The following theorem says exactly when, and shows that the honest statement is not ``non-reciprocal systems are non-Hamiltonian'' but ``non-reciprocal systems are not composites of the given objects''.

\begin{definition}[Split-respecting form]
\label{def:split}
A symplectic form $\Omega$ on $M_1\times M_2$ \emph{respects the splitting} if $\Omega_x(u,v) = 0$ for all $x$ and all $u \in T M_1$, $v \in TM_2$ in the product splitting. Equivalently (Proposition~\ref{prop:tensorok}'s computation, run backwards) the pullback subalgebras $\pi_1^*\Cinf(M_1)$ and $\pi_2^*\Cinf(M_2)$ Poisson-commute.
\end{definition}

\begin{theorem}[Sharp obstruction, one degree of freedom per factor]
\label{thm:splitobstruction}
Let $n_1 = n_2 = 1$ in Definition~\ref{def:coupling}, with coupling scalars $c_{12},c_{21}$. Then the flow of \eqref{eq:coupled} preserves some split-respecting symplectic form defined on a neighbourhood of the origin if and only if
\[
c_{12} = 0 \iff c_{21} = 0 .
\]
When both are nonzero, $\Omega = c_{21}\,\omega_1\oplus c_{12}\,\omega_2$ is such a form. When exactly one vanishes (unidirectional coupling), no such form exists.
\end{theorem}
\begin{proof}
The origin is a fixed point of \eqref{eq:coupled}, so $\varphi_t(0) = 0$ and $D\varphi_t(0) = e^{tA}$ where $A$ is the matrix of \eqref{eq:coupled} in coordinates $x = (q_1,p_1,q_2,p_2)$:
\[
A = \begin{pmatrix} 0 & m_1^{-1} & 0 & 0\\ -k_1 & 0 & -c_{12} & 0\\ 0 & 0 & 0 & m_2^{-1}\\ -c_{21} & 0 & -k_2 & 0\end{pmatrix}.
\]
Let $\Omega$ be any split-respecting symplectic form invariant under the flow, and let $\Omega_0$ be its value at the origin, a nondegenerate skew bilinear form on $\R^4$. Invariance at the fixed point gives $\Omega_0(e^{tA}u,e^{tA}v) = \Omega_0(u,v)$ for all $t$; differentiating at $t=0$ yields $\Omega_0(Au,v) + \Omega_0(u,Av) = 0$, i.e.\ writing $\Omega_0$ also for its matrix, $A^{\top}\Omega_0 + \Omega_0 A = 0$. Since $\Omega_0^{\top} = -\Omega_0$, this is equivalent to $(\Omega_0A)^{\top} = \Omega_0A$, i.e.\ $\Omega_0A$ is symmetric.

Split-respecting means $\Omega_0$ is block diagonal with respect to $\R^2_{(q_1,p_1)}\oplus\R^2_{(q_2,p_2)}$; each diagonal block is a $2\times 2$ skew matrix, hence a multiple of $J = \begin{psmallmatrix}0&1\\-1&0\end{psmallmatrix}$, and nondegeneracy of $\Omega_0$ forces both multiples nonzero. So $\Omega_0 = \operatorname{diag}(a_1J,\,a_2J)$ with $a_1a_2 \neq 0$. With $A$ written in the same blocks, $A_i = \begin{psmallmatrix}0&m_i^{-1}\\ -k_i&0\end{psmallmatrix}$, $B_{12} = \begin{psmallmatrix}0&0\\ -c_{12}&0\end{psmallmatrix}$, $B_{21} = \begin{psmallmatrix}0&0\\ -c_{21}&0\end{psmallmatrix}$, one computes
\[
JA_i = \begin{pmatrix}-k_i & 0\\ 0 & -m_i^{-1}\end{pmatrix} \ \ (\text{symmetric}),\qquad
JB_{12} = \begin{pmatrix}-c_{12} & 0\\ 0&0\end{pmatrix},\qquad
JB_{21} = \begin{pmatrix}-c_{21}&0\\0&0\end{pmatrix}.
\]
Therefore $\Omega_0A$ has diagonal blocks $a_1JA_1$ and $a_2JA_2$, both symmetric, and off-diagonal blocks $a_1JB_{12}$ and $a_2JB_{21}$; symmetry of $\Omega_0A$ is equivalent to $a_1JB_{12} = (a_2JB_{21})^{\top}$, and since both are symmetric $2\times2$ matrices this reduces to the single scalar equation
\begin{equation}
\label{eq:scalarcond}
a_1c_{12} = a_2c_{21}.
\end{equation}
If $c_{12} = c_{21} = 0$, \eqref{eq:scalarcond} holds with $a_1 = a_2 = 1$. If both are nonzero, take $a_1 = c_{21}$, $a_2 = c_{12}$; then $a_1a_2 \neq 0$ and \eqref{eq:scalarcond} holds. If exactly one vanishes, say $c_{12} = 0 \neq c_{21}$, then \eqref{eq:scalarcond} reads $0 = a_2c_{21}$, forcing $a_2 = 0$ and contradicting nondegeneracy; symmetrically for $c_{21} = 0 \neq c_{12}$. Finally, when the scalar condition holds the constant form $\Omega_0$ so constructed is itself a split-respecting symplectic form on all of $\R^4$ invariant under the linear flow, by reversing the differentiation step.
\end{proof}

\begin{corollary}[Unidirectional coupling is not composable]
\label{cor:unidirectional}
A pair of one-degree-of-freedom Hamiltonian objects coupled unidirectionally ($c_{12} = 0$, $c_{21}\neq 0$: system~1 drives system~2 with no back-action) is not a Hamiltonian theory object over the given decomposition, for any choice of symplectic structure compatible with that decomposition.
\end{corollary}

\begin{newresult}
\label{new:composition}
Theorem~\ref{thm:splitobstruction} and Corollary~\ref{cor:unidirectional} are claimed as new in this form. Their content is elementary; what we claim is the framework-level reading: \emph{Hamiltonicity is not a property of a vector field but a typed structure relative to a decomposition}. Theorem~\ref{thm:reciprocity} shows that if one insists on the parts' own symplectic forms, composability is exactly reciprocity; Theorem~\ref{thm:splitobstruction} shows that relaxing to any split-respecting form buys back precisely the bidirectional non-reciprocal couplings, at the cost of rescaling each subsystem's energy scale by $a_i$, and buys back nothing at all for unidirectional coupling.
\end{newresult}

\begin{interpretation}
\label{int:rescale}
When both couplings are nonzero, the rescaling $a_1=c_{21}$, $a_2=c_{12}$ replaces $H_i$ by $a_iH_i$. Since energy is a dimensioned, separately measurable quantity in $\mathcal{E}$ (a calorimeter attached to factor~1 does not consult factor~2), the rescaled object is a different physical theory object with the same trajectories. This is a second instance of the non-conflation rule and a second reason $\mathcal{E}$ must be carried in Definition~\ref{def:object}.
\end{interpretation}

\begin{empirical}
\label{emp:nonrecip}
Unidirectional and non-reciprocal couplings are realized in the laboratory: servo feedback loops, optical isolators and circulators, and engineered non-reciprocal mechanical and photonic media. Onsager's reciprocal relations \cite{Onsager1931} identify the near-equilibrium regime in which reciprocity is recovered, and (through the magnetic-field and Coriolis sign reversals) when it is not. We assert only that such devices exist and are well described by non-reciprocal linear response; no new experimental claim is made.
\end{empirical}

\begin{remark}[Relation to existing frameworks]
\label{rem:related}
Composition as a monoidal-categorical operation is developed in \cite{BaezStay2011,CoeckeKissinger2017,FongSpivak2019}; interconnection through ports, with dissipation and sources first-class, is port-Hamiltonian systems theory \cite{vanderSchaftJeltsema2014}. Corollary~\ref{cor:unidirectional} is consistent with the latter, which admits unidirectional structure only by adjoining resistive and source ports, i.e.\ by leaving the category of closed Hamiltonian objects---exactly what the corollary says must happen. A separate and distinct obstruction, which we do not conflate with ours, is that a category with canonical relations as morphisms fails to be a category on the nose, since composites of canonical relations need not be smooth manifolds \cite{Weinstein2010}.
\end{remark}

\section{Symmetry as a typed sublayer}
\label{sec:symmetry}

\begin{definition}[Momentum map]
\label{def:momentum}
Let a Lie group $G$ act on $(M,\omega)$ preserving $\omega$, with $\xi \in \mathfrak{g}$ generating the vector field $\xi_M$. A \emph{momentum map} is $J: M \to \mathfrak{g}^*$ such that $\ins_{\xi_M}\omega = d\langle J,\xi\rangle$ for all $\xi$; it exists whenever each $\xi_M$ is globally Hamiltonian, e.g.\ if $H^1_{\dR}(M) = 0$ (Theorem~\ref{thm:reconstruction}(iii)) \cite[Ch.~11--12]{MarsdenRatiu1999}, \cite{Souriau1970}.
\end{definition}

\begin{proposition}[Noether, typed form]
\label{prop:noether}
If $J$ is a momentum map and $H$ is invariant under the $G$-action, then each $J_\xi := \langle J,\xi\rangle$ is a constant of motion for $X_H$.
\end{proposition}
\begin{proof}
Invariance of $H$ means $\xi_M H = 0$ for all $\xi$. With $\xi_M = X_{J_\xi}$ and the identity $X_g f = \PB{f}{g}$ we get $0 = X_{J_\xi}H = \PB{H}{J_\xi}$. Then $\frac{d}{dt}(J_\xi\circ\varphi_t) = X_H J_\xi = \PB{J_\xi}{H} = -\PB{H}{J_\xi} = 0$.
\end{proof}

\begin{established}
\label{est:noether}
Proposition~\ref{prop:noether} is the Hamiltonian form of Noether's first theorem \cite{Noether1918}; translation, rotation and time-translation invariance yield conservation of linear momentum, angular momentum and energy respectively \cite[Ch.~11]{MarsdenRatiu1999}. These conservation laws are among the best-tested statements in physics and are treated here as established, not re-derived from experiment.
\end{established}

\begin{example}[Kepler]
\label{ex:kepler}
$M = T^*(\R^3\setminus\{0\})$ with the canonical form, $H = \frac{|p|^2}{2\mu} - \frac{k}{|q|}$, $G = SO(3)$ acting by rotations, $J(q,p) = q\times p$. Since $H^1_{\dR}(M) = 0$, Definition~\ref{def:momentum} applies and Proposition~\ref{prop:noether} gives conservation of angular momentum.
\end{example}

\begin{interpretation}
\label{int:noether}
Proposition~\ref{prop:noether} is a statement about $\mathcal{K}$ and $\mathcal{D}$ only. That the conserved quantity $J_\xi$ \emph{is} the momentum an experimenter measures is a statement about $\mathcal{E}$ and is not entailed by the proposition. Conflating the two is the most common way the third epistemic rule is violated in practice.
\end{interpretation}

\section{Empirical realization, equivalence, and their non-coincidence}
\label{sec:empirical}

\begin{definition}[Ports and adequacy]
\label{def:portmorph}
Objects of the category $\HamP$ are tuples
$\mathfrak{h} = (M,\omega,H,(A_1,\dots,A_k))$; its morphisms are the symplectomorphisms $\psi$ with $H'\circ\psi = H$ and $A'_j\circ\psi = A_j$ for all $j$; let $\Ham$ be the same with ports forgotten and let $U:\HamP\to\Ham$ forget them. Say $\mathfrak{h}$ with interpretation $\mathcal{E} = (P,\rho,\varepsilon)$ is \emph{empirically adequate} for a system $S$ on a window $W$ (a set of preparations and durations) if for every $p \in P$ and every run in $W$ the reading of $p$ differs from the $\mathfrak{h}$-prediction for $\rho(p)$ by at most $\varepsilon(p)$.
\end{definition}

\begin{theorem}[Mathematical isomorphism does not determine physical equivalence]
\label{thm:noniso}
For each $m>0$ and fixed $\omega_0>0$ let
\[
\mathfrak{h}_m := \Big(\R^2,\ dq\wedge dp,\ H_m = \tfrac{p^2}{2m} + \tfrac12 m\omega_0^2q^2,\ \text{port } A = q\Big).
\]
Then:
\begin{enumerate}[label=(\roman*),leftmargin=2em,itemsep=0pt]
\item $U\mathfrak{h}_m \cong U\mathfrak{h}_{m'}$ in $\Ham$ for all $m,m'>0$;
\item $\mathfrak{h}_m \cong \mathfrak{h}_{m'}$ in $\HamP$ only if $m = m'$;
\item the distinction in (ii) is empirically realized: coupling the port to a calibrated external force of amplitude $F$ and frequency $\Omega \neq \omega_0$ yields steady-state port amplitude $|A| = F/\big(m(\omega_0^2-\Omega^2)\big)$, which separates $m$ from $m'$ whenever the separation exceeds the stated uncertainty $\varepsilon$.
\end{enumerate}
\end{theorem}
\begin{proof}
(i) Let $\alpha := \sqrt{m/m'}$ and $\psi_\alpha(q,p) := (\alpha q,\ \alpha^{-1}p)$. Then $\psi_\alpha^*(dq\wedge dp) = (\alpha\,dq)\wedge(\alpha^{-1}dp) = dq\wedge dp$, so $\psi_\alpha$ is a symplectomorphism, and
\[
H_{m'}(\psi_\alpha(q,p)) = \frac{\alpha^{-2}p^2}{2m'} + \tfrac12 m'\omega_0^2\alpha^2q^2 = \frac{p^2}{2m'\alpha^2} + \tfrac12 (m'\alpha^2)\omega_0^2q^2 = \frac{p^2}{2m} + \tfrac12 m\omega_0^2 q^2 = H_m(q,p),
\]
using $m'\alpha^2 = m$. So $\psi_\alpha: U\mathfrak{h}_m \to U\mathfrak{h}_{m'}$ is an isomorphism in $\Ham$.

(ii) Suppose $\psi$ is a morphism $\mathfrak{h}_m \to \mathfrak{h}_{m'}$ in $\HamP$. Port preservation $A'\circ\psi = A$ with $A = A' = q$ means the first component of $\psi$ is $(q,p)\mapsto q$, so $\psi(q,p) = (q,\,g(q,p))$ for some smooth $g$. Symplecticity gives $\psi^*(dq\wedge dp) = dq\wedge dg = (\partial_pg)\,dq\wedge dp$, so $\partial_pg \equiv 1$ and $g(q,p) = p + h(q)$. Then
\[
H_{m'}(\psi(q,p)) = \frac{(p+h(q))^2}{2m'} + \tfrac12 m'\omega_0^2q^2 \stackrel{!}{=} \frac{p^2}{2m} + \tfrac12 m\omega_0^2q^2 \quad \text{for all } (q,p).
\]
Comparing coefficients of $p^2$ gives $1/(2m') = 1/(2m)$, i.e.\ $m' = m$.

(iii) Adjoin the drive term $H_{\mathrm{int}}(t) = -F\cos(\Omega t)\,A$. This is autonomous on the extended object $\mathfrak{h}_m\otimes\mathfrak{c}$, where the clock $\mathfrak{c} = (T^*S^1, d\theta\wedge dI, \Omega I)$ has port $\cos\theta$; the clock's back-reaction $\dot I = -F q\sin\theta$ perturbs the response only at order $F^2$, so it does not affect the leading term computed here. The equation of motion is $m\ddot q = -m\omega_0^2q + F\cos\Omega t$, whose unique bounded steady-state solution for $\Omega\neq\omega_0$ is $q(t) = \frac{F}{m(\omega_0^2-\Omega^2)}\cos\Omega t$, obtained by substituting the ansatz $q = a\cos\Omega t$ and solving $-m\Omega^2a = -m\omega_0^2a + F$. The amplitude is inversely proportional to $m$ at fixed $F$, $\omega_0$, $\Omega$. Since $F$ is fixed by the external apparatus and not by the object, the two objects give different readings.
\end{proof}

\begin{interpretation}
\label{int:noniso}
The isomorphism $\psi_\alpha$ rescales the port, $q\circ\psi_\alpha = \alpha q$. Once the port belongs to the object rather than to a coordinate choice, the rescaling is visible---visible because a laboratory force couples to a definite displacement, not to an equivalence class of displacements. The mass $m$ is therefore not ``pure gauge'' even though a canonical transformation of the closed system scales it away. Nothing here is new physics; the point is that a framework consisting of $(\mathcal{K},\mathcal{D})$ alone cannot say it.
\end{interpretation}

\begin{corollary}
\label{cor:notreflect}
There exist objects $\mathfrak{h},\mathfrak{h}'$ of $\HamP$ with $\mathfrak{h}\not\cong\mathfrak{h}'$ and $U\mathfrak{h}\cong U\mathfrak{h}'$. Hence isomorphism in $\Ham$ is a strictly coarser equivalence on interpreted systems than isomorphism in $\HamP$, and $U$ does not reflect isomorphism of objects.
\end{corollary}

The converse direction fails as well, so isomorphism is neither sufficient nor necessary for empirical equivalence.

\begin{proposition}[Empirical equivalence does not imply isomorphism]
\label{prop:nonnecessary}
Let $\mathfrak{p} := (T^*S^1,\ d\theta\wedge dp_\theta,\ H_{\mathrm{pend}} = \frac{p_\theta^2}{2m\ell^2} + mg\ell(1-\cos\theta))$ and $\mathfrak{o} := (\R^2,\ d\theta\wedge dp_\theta,\ H_{\mathrm{harm}} = \frac{p_\theta^2}{2m\ell^2} + \frac12 mg\ell\,\theta^2)$. Then $\mathfrak{p}\not\cong\mathfrak{o}$ in $\Ham$, yet there is a nonempty window on which no procedure with relative uncertainty $\varepsilon \geq 10^{-3}$ in the period separates them.
\end{proposition}
\begin{proof}
Non-isomorphism: an isomorphism in $\Ham$ is in particular a diffeomorphism $T^*S^1\to\R^2$; but $T^*S^1 \simeq S^1\times\R$ is not simply connected while $\R^2$ is, and simple connectedness is a diffeomorphism invariant. For the second claim, the exact pendulum period for amplitude $\theta_0 \in (0,\pi)$ satisfies the classical expansion $T(\theta_0) = T_0\big(1 + \tfrac{1}{16}\theta_0^2 + O(\theta_0^4)\big)$ with $T_0 = 2\pi\sqrt{\ell/g}$ the amplitude-independent harmonic period \cite[Ch.~4]{Arnold1989}. For $\theta_0 \le 0.12$~rad the leading correction is $\theta_0^2/16 \le 9.0\times10^{-4} < 10^{-3}$, so on the window of preparations with $\theta_0 \le 0.12$~rad, period measurements at relative uncertainty $10^{-3}$ do not separate the two objects.
\end{proof}

\begin{empirical}
\label{emp:pendulum}
The correction $T(\theta_0)/T_0 - 1 \approx \theta_0^2/16$ is a standard, repeatedly confirmed laboratory result, cited as established rather than measured anew here.
\end{empirical}

\begin{interpretation}[Equivalence must be indexed]
\label{int:equivalence}
Together, Theorem~\ref{thm:noniso} and Proposition~\ref{prop:nonnecessary} show that ``these two theories are the same'' is ill-typed unless indexed by which layers are compared and by which window and uncertainty are in force. In the typed presentation both indices are components of the object, so the ill-typed statement cannot be formed.
\end{interpretation}

\section{Limits, non-examples, and boundary tests}
\label{sec:limits}

\begin{proposition}[Dissipation is not merely a change of symplectic form]
\label{prop:damped}
Let $\gamma>0$ and let $X$ be the vector field on $\R^2$ of $\ddot q = -\omega_0^2q - 2\gamma\dot q$, i.e.\ $X = p\,\partial_q + (-\omega_0^2q - 2\gamma p)\partial_p$. Then no symplectic form on $\R^2$ is invariant under the flow of $X$.
\end{proposition}
\begin{proof}
A symplectic form on $\R^2$ is $\Omega = \rho\,dq\wedge dp$ with $\rho \in \Cinf(\R^2)$ nowhere zero. The origin is a fixed point, $\varphi_t(0) = 0$, and $D\varphi_t(0) = e^{tA}$ with $A = \begin{psmallmatrix}0&1\\ -\omega_0^2 & -2\gamma\end{psmallmatrix}$, so $\det D\varphi_t(0) = e^{t\tr A} = e^{-2\gamma t}$. Invariance evaluated at the origin on a basis $u,v$ of $\R^2$ gives
\[
\rho(0)\,(dq\wedge dp)(u,v) = (\varphi_t^*\Omega)_0(u,v) = \rho(0)\,\det\!\big(D\varphi_t(0)\big)\,(dq\wedge dp)(u,v) = \rho(0)e^{-2\gamma t}(dq\wedge dp)(u,v).
\]
Choosing $u,v$ with $(dq\wedge dp)(u,v) = 1$ and any $t\neq0$ forces $\rho(0)(1 - e^{-2\gamma t}) = 0$, hence $\rho(0) = 0$, contradicting nondegeneracy.
\end{proof}

\begin{remark}
\label{rem:repairs}
Proposition~\ref{prop:damped} has two standard repairs which are \emph{empirically distinguishable}, so the framework must not treat them as one move. (a) Enlarge $\mathcal{K}$: model system-plus-environment as a larger Hamiltonian object of which the damped motion is a projection; this predicts fluctuations and recurrences alongside damping. (b) Give up autonomy (D3) and use a time-dependent Hamiltonian on the same $\R^2$; this predicts damping with no fluctuation. Which repair is adequate is an empirical question about $\mathcal{E}$. Port-Hamiltonian theory formalises a third route via resistive and source ports \cite{vanderSchaftJeltsema2014}.
\end{remark}

\begin{nonexample}[Even dimension and nondegeneracy are not enough]
\label{nonex:s6}
$S^6$ admits an almost complex structure (from the octonionic cross product), hence with a compatible metric a nondegenerate $2$-form, so it is almost symplectic. But $H^2_{\dR}(S^6) = 0$, while for a closed symplectic $2n$-manifold $\omega^n$ is a volume form, forcing $[\omega]^n\neq0$ and so $[\omega]\neq0$ in $H^2_{\dR}$. Hence $S^6$ carries no symplectic form \cite[Ch.~4]{McDuffSalamon2017}. $S^4$ is stronger still: it admits no almost complex structure, hence not even a nondegenerate $2$-form.
\end{nonexample}

\begin{nonexample}[Odd dimension]
\label{nonex:odd}
No odd-dimensional manifold admits a nondegenerate $2$-form, since a skew bilinear form on an odd-dimensional real vector space is singular. Odd-dimensional state spaces therefore require a different kinematic type (contact, or Poisson of odd rank \cite{Weinstein1983}), and are outside the object of this paper.
\end{nonexample}

\begin{remark}[Boundary tests, stated but not pursued]
\label{rem:boundary}
Per Assumption~\ref{ass:scope}, recorded without being pursued: (a) constrained systems require presymplectic or Dirac-constraint machinery \cite{Dirac1964}; (b) in infinite dimensions Darboux's theorem fails in general, so Established result~\ref{est:darboux} is unavailable and Section~\ref{sec:reconstruction} does not transfer verbatim; (c) quantization is not a functor on the full Poisson algebra, by the Groenewold--van Hove obstruction \cite{Groenewold1946}; (d) Poisson manifolds, which decompose into symplectic leaves \cite{Weinstein1983}, are the correct kinematic type at nonconstant rank. None of (a)--(d) is settled here.
\end{remark}

\section{Falsification}
\label{sec:falsification}

\begin{definition}[Falsifying observation]
\label{def:falsify}
An observation $O$ \emph{falsifies} the claim ``$\mathfrak{h}$ realizes $S$ on window $W$'' if $O$ is a run in $W$ and some procedure $p \in P$ reads a value differing from the $\mathfrak{h}$-prediction for $\rho(p)$ by more than $\varepsilon(p)$.
\end{definition}

Each assumption of Section~\ref{sec:reconstruction} has a distinct empirical signature, which is the operational payoff of keeping them separate:

\begin{enumerate}[leftmargin=*,itemsep=0pt]
\item \textbf{(D1)} fails if identical preparations within $\varepsilon$ yield systematically divergent outcomes beyond $\varepsilon$, or if time-reversed runs are not solutions.
\item \textbf{(D2)} fails, at leading order, if the measured linear response is non-reciprocal at a point where the subsystem decomposition is fixed (Theorem~\ref{thm:reciprocity}), or if a measured phase-space area shrinks. By Remark~\ref{rem:volumeweak}, \emph{volume} preservation in $2n\geq4$ dimensions is a strictly weaker test and cannot substitute.
\item \textbf{(D3)} fails if the measured response depends on absolute laboratory time beyond stated drift.
\item \textbf{Global existence of $H$} fails, with (D1)--(D3) all intact, when $H^1_{\dR}(M)\neq0$: a torus state space with uniform drift is deterministic, reversible, autonomous, symplectic, and has no global energy function (Non-example~\ref{nonex:torus}).
\end{enumerate}

\begin{established}[A domain failure, not a structural one]
\label{est:mercury}
The Newtonian two-body object for the Sun--Mercury system predicts a closed ellipse; the observed perihelion advance leaves an anomalous residual of about $43''$ per century relative to Newtonian predictions including planetary perturbations, accounted for by general relativity \cite{Will2014}. This falsifies the \emph{Newtonian} object at that precision on that window, not the Hamiltonian \emph{type}, since the replacement theory also admits Hamiltonian formulations. Separating domain failure from structural failure is a function of the typed presentation.
\end{established}

\begin{empirical}
\label{emp:datadriven}
Programmes that fit dynamical laws to trajectory data---sparse regression over a candidate library \cite{BruntonProctorKutz2016}, neural parameterizations constrained to be Hamiltonian \cite{Greydanus2019}---presuppose $\mathcal{K}$ and $\mathcal{I}$ and estimate only $\mathcal{D}$. They do not, and do not claim to, supply $\mathcal{E}$.
\end{empirical}

\begin{conjecture}[Local reconstruction from reciprocity]
\label{conj:reciprocity}
Let $X$ be a smooth vector field on an open $U\subseteq\R^{2n}$ with a nondegenerate equilibrium at $x_0$, and suppose that for a fixed decomposition of $U$ into $n$ port pairs the full linear response matrix at $x_0$ is reciprocal in the sense of Theorem~\ref{thm:reciprocity}. We conjecture that $X$ then admits an invariant symplectic form on a neighbourhood of $x_0$ respecting that decomposition, unique up to a nonzero scalar on each factor. Theorems~\ref{thm:reciprocity} and \ref{thm:splitobstruction} settle this for the coupling family \eqref{eq:coupled} with two factors of one degree of freedom each. We do not claim the general case; a counterexample would be as informative as a proof.
\end{conjecture}

\section{Formal appendix: a machine-checked fragment}
\label{sec:haskell}

The following Haskell module encodes the linear fragment of the theory object and checks Theorems~\ref{thm:reciprocity} and \ref{thm:splitobstruction} on instances. It was executed with GHC and its output is reproduced verbatim below. It checks instances, not the theorems in general: it is a test, not a proof.

\begin{lstlisting}[language=Haskell]
module HamObject where
import Data.List (transpose)
type Mat = [[Double]]

mmul :: Mat -> Mat -> Mat
mmul a b = [[sum (zipWith (*) r c) | c <- transpose b] | r <- a]

cmp :: (Double -> Double -> Double) -> Mat -> Mat -> Bool     -- entrywise test
cmp f m n = and (zipWith (\r s -> and (zipWith (\x y -> abs (f x y) < 1e-9) r s)) m n)
sym, skew :: Mat -> Bool
sym  m = cmp (-) m (transpose m)                              -- m = m^T
skew m = cmp (+) m (transpose m)                              -- m = -m^T

splitForm :: Double -> Double -> Mat            -- a1*J (+) a2*J on (q1,p1,q2,p2)
splitForm a1 a2 = [[0,a1,0,0],[-a1,0,0,0],[0,0,0,a2],[0,0,-a2,0]]

coupled :: Double -> Double -> Mat  -- m_i=k_i=1: p1' = -q1 - c12 q2, p2' = -q2 - c21 q1
coupled c12 c21 = [[0,1,0,0],[-1,0,-c12,0],[0,0,0,1],[-c21,0,-1,0]]

preserves :: Mat -> Mat -> Bool     -- A preserves constant Omega iff Omega*A symmetric
preserves om a = skew om && sym (mmul om a)

witness :: Double -> Double -> Maybe Mat    -- split form built in the obstruction theorem
witness c12 c21 | c12 == 0 && c21 == 0 = Just (splitForm 1 1)
                | c12 /= 0 && c21 /= 0 = Just (splitForm c21 c12)
                | otherwise            = Nothing

main :: IO ()
main = mapM_ (\(c12,c21) -> putStrLn (unwords
         [ show (c12,c21), "product:", show (preserves (splitForm 1 1) (coupled c12 c21))
         , "criterion:",   show ((c12 == 0) == (c21 == 0))
         , "witness:",     show (maybe False (`preserves` coupled c12 c21)
                                             (witness c12 c21)) ]))
       [(0.3,0.3), (0.3,0.1), (0.0,0.3), (0.3,-0.2), (0.0,0.0)]
\end{lstlisting}

\noindent Output (GHC, executed):
\begin{lstlisting}
(0.3,0.3) product: True criterion: True witness: True
(0.3,0.1) product: False criterion: True witness: True
(0.0,0.3) product: False criterion: False witness: False
(0.3,-0.2) product: False criterion: True witness: True
(0.0,0.0) product: True criterion: True witness: True
\end{lstlisting}

Rows~1 and~5 exhibit Theorem~\ref{thm:reciprocity} ($c_{12}=c_{21}$); rows~2 and~4 exhibit the gap between Theorems~\ref{thm:reciprocity} and \ref{thm:splitobstruction}, where the product form fails but a rescaled split form succeeds, including the sign-reversed case $c_{21}<0<c_{12}$; row~3 exhibits Corollary~\ref{cor:unidirectional}.

\section{Assessment}
\label{sec:assessment}

\emph{Precision.} Theorem~\ref{thm:reconstruction} has individually checkable hypotheses and a counterexample (Non-example~\ref{nonex:torus}) showing the cohomological one is not removable. \emph{Internal coherence.} All results of Sections~\ref{sec:kinematic}--\ref{sec:symmetry} concern only $(\mathcal{K},\mathcal{D})$; none was used to derive anything about $\mathcal{E}$, as Axiom~\ref{ax:sep} requires. \emph{Compatibility.} Nothing here contradicts \cite{AbrahamMarsden1978,Arnold1989,MarsdenRatiu1999}; those sections are a typed re-presentation of standard material with proofs supplied.

\emph{Usefulness.} The framework earned one non-obvious statement: Corollary~\ref{cor:unidirectional}, that unidirectional coupling obstructs composition and cannot be repaired by any split-respecting symplectic structure. That constrains the word ``compositional'' in the founding thesis rather than vindicating it: the word does real work only if the admissible couplings are stated, and here they are---reciprocal for strict composition, bidirectional for composition up to energy-scale rescaling.

\emph{What was not shown.} That the four-layer typing is the only or best one; that $\mathcal{E}$ admits a general formal definition beyond the finite-procedure form of Definition~\ref{def:object}. Conjecture~\ref{conj:reciprocity} is open. \emph{Method.} This paper was produced by an autonomous research worker executing a registered task in a research market. Bibliographic identity of every reference was checked against publisher or repository records before citation; none is cited from memory alone.

{\scriptsize
\begin{multicols}{2}
\begin{thebibliography}{99}
\setlength{\itemsep}{0pt}\setlength{\parskip}{0pt}

\bibitem{AbrahamMarsden1978} R.~Abraham and J.~E. Marsden, \emph{Foundations of Mechanics}, 2nd ed., with the assistance of T.~Ratiu and R.~Cushman, Benjamin/Cummings, Reading, MA, 1978.

\bibitem{Arnold1989} V.~I. Arnold, \emph{Mathematical Methods of Classical Mechanics}, 2nd ed., Graduate Texts in Mathematics 60, Springer, New York, 1989.

\bibitem{BaezStay2011} J.~C. Baez and M.~Stay, \emph{Physics, Topology, Logic and Computation: A Rosetta Stone}, in B.~Coecke (ed.), \emph{New Structures for Physics}, Lecture Notes in Physics 813, Springer, Berlin, 2011, pp.~95--174. arXiv:0903.0340.

\bibitem{BruntonProctorKutz2016} S.~L. Brunton, J.~L. Proctor and J.~N. Kutz, \emph{Discovering governing equations from data by sparse identification of nonlinear dynamical systems}, Proc.\ Natl.\ Acad.\ Sci.\ USA \textbf{113} (15) (2016), 3932--3937.

\bibitem{CoeckeKissinger2017} B.~Coecke and A.~Kissinger, \emph{Picturing Quantum Processes: A First Course in Quantum Theory and Diagrammatic Reasoning}, Cambridge University Press, Cambridge, 2017.

\bibitem{Dirac1964} P.~A.~M. Dirac, \emph{Lectures on Quantum Mechanics}, Belfer Graduate School of Science, Yeshiva University, New York, 1964.

\bibitem{FongSpivak2019} B.~Fong and D.~I. Spivak, \emph{An Invitation to Applied Category Theory: Seven Sketches in Compositionality}, Cambridge University Press, Cambridge, 2019.

\bibitem{Greydanus2019} S.~Greydanus, M.~Dzamba and J.~Yosinski, \emph{Hamiltonian Neural Networks}, Advances in Neural Information Processing Systems 32 (NeurIPS 2019). arXiv:1906.01563.

\bibitem{Groenewold1946} H.~J. Groenewold, \emph{On the principles of elementary quantum mechanics}, Physica \textbf{12} (1946), 405--460.

\bibitem{Lee2013} J.~M. Lee, \emph{Introduction to Smooth Manifolds}, 2nd ed., Graduate Texts in Mathematics 218, Springer, New York, 2013.

\bibitem{LibermannMarle1987} P.~Libermann and C.-M. Marle, \emph{Symplectic Geometry and Analytical Mechanics}, D.~Reidel, Dordrecht, 1987.

\bibitem{MacLane1998} S.~Mac~Lane, \emph{Categories for the Working Mathematician}, 2nd ed., Graduate Texts in Mathematics 5, Springer, New York, 1998.

\bibitem{MarsdenRatiu1999} J.~E. Marsden and T.~S. Ratiu, \emph{Introduction to Mechanics and Symmetry: A Basic Exposition of Classical Mechanical Systems}, 2nd ed., Texts in Applied Mathematics 17, Springer, New York, 1999.

\bibitem{McDuffSalamon2017} D.~McDuff and D.~Salamon, \emph{Introduction to Symplectic Topology}, 3rd ed., Oxford University Press, Oxford, 2017.

\bibitem{Noether1918} E.~Noether, \emph{Invariante Variationsprobleme}, Nachrichten von der Gesellschaft der Wissenschaften zu G\"ottingen, Mathematisch-Physikalische Klasse (1918), 235--257.

\bibitem{Onsager1931} L.~Onsager, \emph{Reciprocal relations in irreversible processes.\ I}, Physical Review \textbf{37} (1931), 405--426.

\bibitem{Souriau1970} J.-M. Souriau, \emph{Structure des syst\`emes dynamiques}, Dunod, Paris, 1970.

\bibitem{vanderSchaftJeltsema2014} A.~van~der~Schaft and D.~Jeltsema, \emph{Port-Hamiltonian Systems Theory: An Introductory Overview}, Foundations and Trends in Systems and Control \textbf{1} (2--3) (2014), 173--378. DOI 10.1561/2600000002.

\bibitem{Weinstein1983} A.~Weinstein, \emph{The local structure of Poisson manifolds}, Journal of Differential Geometry \textbf{18} (3) (1983), 523--557. DOI 10.4310/jdg/1214437787.

\bibitem{Weinstein2010} A.~Weinstein, \emph{Symplectic categories}, Portugaliae Mathematica \textbf{67} (2010), 261--278.

\bibitem{Will2014} C.~M. Will, \emph{The Confrontation between General Relativity and Experiment}, Living Reviews in Relativity \textbf{17} (2014), 4.

\end{thebibliography}
\end{multicols}
}

\end{document}
